\documentclass[runningheads]{llncs}

\usepackage[spanish]{babel}
\usepackage{url}

\renewcommand{\andname}{y}
\renewcommand{\lastandname}{ y}
\renewcommand{\keywordname}{{\bfseries Palabras clave:}}

\begin{document}

\title{Transpilador-Model-SQL: un transpilador de modelos relacionales a SQL}
\author{Alan Zschach\inst{1} \and Agustín Mango\inst{1} \and
Zahira Alves\inst{1} \and Erick Ezequiel Kusczak\inst{1}}
\authorrunning{A. Zschach y cols.}
\institute{Universidad Gastón Dachary\\Compiladores 2026}
\maketitle

\begin{abstract}
Transpilador-Model-SQL es un transpilador de un lenguaje específico de dominio (DSL) declarativo
para describir esquemas relacionales. A partir de archivos de texto \texttt{.model},
la herramienta verifica la estructura y las restricciones del esquema y genera
sentencias de definición de datos (DDL) para PostgreSQL. El análisis léxico se
implementa en C con Flex y el análisis sintáctico con Bison. El árbol de sintaxis
abstracta resultante se valida semánticamente en dos pasadas y, si no hay errores,
se traduce a SQL. El proyecto incluye diagnósticos con ubicación, una interfaz de
línea de comandos, un Makefile que genera todos los componentes y escenarios de
integración válidos e inválidos ejecutables con un único comando.

\keywords{compiladores \and lenguajes específicos de dominio \and Flex \and Bison \and SQL}
\end{abstract}

\section{Introducción}

La definición manual de esquemas relacionales mediante DDL requiere repetir
estructuras y mantener consistentes tipos, claves y referencias. Transpilador-Model-SQL aborda
ese problema con un lenguaje pequeño que expresa modelos, atributos y restricciones,
y traduce las declaraciones a SQL. La herramienta automatiza la generación de
\texttt{CREATE TABLE} y detecta errores antes de producir el resultado.

El trabajo aplica las etapas habituales de un compilador: análisis léxico,
análisis sintáctico, construcción de una representación intermedia, análisis
semántico y generación de código \cite{aho}. La consigna requiere el uso de
generadores de analizadores léxicos y sintácticos, por lo que la etapa de análisis
se construye con Flex y Bison \cite{flex,bison}.

\section{Antecedentes}

Un compilador separa el reconocimiento de lexemas, la verificación de la
estructura gramatical y la comprobación de restricciones que dependen del
significado del programa. La separación permite reportar errores en la etapa
correspondiente y evitar que la etapa de generación reciba un esquema inválido
\cite{aho}.

Flex genera un analizador léxico a partir de reglas basadas en expresiones
regulares; Bison genera un analizador sintáctico a partir de una gramática libre
de contexto \cite{flex,bison}. Transpilador-Model-SQL utiliza esas herramientas
para convertir los elementos léxicos en un árbol de sintaxis abstracta. La etapa
de generación produce DDL para PostgreSQL,
cuyas restricciones determinan, entre otros aspectos, que las columnas
referenciadas por claves foráneas sean claves primarias o únicas
\cite{postgresql}.

\section{Lenguaje fuente}

El lenguaje fuente es declarativo. Una declaración comienza con la palabra
reservada \texttt{model}, seguida del nombre del modelo y un bloque entre llaves.
Dentro del bloque, cada atributo declara un nombre, un tipo y cero o más
modificadores. Los tipos disponibles son \texttt{int}, \texttt{string},
\texttt{float}, \texttt{boolean} y \texttt{date}. Los modificadores son
\texttt{pk}, \texttt{unique}, \texttt{not\_null}, \texttt{auto\_increment} y
\texttt{fk(Modelo.atributo)}.

Los espacios, tabulaciones y saltos de línea separan lexemas, pero no delimitan
atributos; los comentarios comienzan con \texttt{//}. La gramática completa se
mantiene en \texttt{docs/gramatica.md}, junto con la correspondencia de tipos
del DSL a PostgreSQL.

\subsection{Reglas semánticas}

Los nombres de modelos son únicos en el esquema y los atributos son únicos dentro
de su modelo. Cada modelo debe tener al menos un atributo y puede declarar como
máximo una clave primaria. El modificador \texttt{auto\_increment} requiere el
tipo \texttt{int}.

Una clave foránea debe referenciar un modelo y un atributo existentes. Los tipos
del atributo local y del atributo referenciado deben coincidir, y el atributo
referenciado debe ser \texttt{pk} o \texttt{unique}. El análisis de dos pasadas
permite resolver referencias a modelos declarados más adelante.

\section{Arquitectura e implementación}

\subsection{Análisis léxico y sintáctico}

El analizador léxico generado con Flex reconoce palabras reservadas, identificadores, signos de
puntuación, espacios y comentarios. Cada token conserva su ubicación de origen.
Los caracteres no reconocidos producen un error léxico con línea y columna.

El analizador sintáctico generado con Bison aplica la gramática del DSL y
construye el árbol de sintaxis abstracta: un esquema contiene modelos; cada modelo
contiene atributos; cada atributo almacena su tipo, posición y modificadores.
Ante una secuencia inválida, informa un error sintáctico con ubicación y el
elemento inesperado.

\subsection{Análisis semántico}

El analizador semántico usa dos pasadas. La primera registra modelos y atributos
en una tabla de símbolos y comprueba duplicados, cantidad de claves primarias,
modelos vacíos y el uso de \texttt{auto\_increment}. La segunda resuelve cada
referencia de clave foránea y comprueba existencia, unicidad y compatibilidad
de tipos. Se informa el primer error y no se genera SQL cuando el esquema no
supera la validación.

\subsection{Generación de código y ejecución}

El generador traduce los tipos a \texttt{INTEGER}, \texttt{VARCHAR(255)},
\texttt{REAL}, \texttt{BOOLEAN} y \texttt{DATE}; \texttt{int} con
\texttt{auto\_increment} se emite como \texttt{SERIAL}. Los modificadores se
traducen a restricciones de clave primaria, unicidad y no nulidad.

Las tablas se emiten en el orden de declaración. Las claves foráneas a tablas
anteriores se incluyen en el \texttt{CREATE TABLE} correspondiente; las
referencias adelantadas se agregan después mediante \texttt{ALTER TABLE}, una
vez creadas todas las tablas. Esta estrategia conserva el orden fuente y permite
generar DDL válido también para relaciones circulares.

El Makefile es el punto de entrada de la compilación: genera el analizador
sintáctico con Bison, el analizador léxico con Flex y el ejecutable
\texttt{build/model2sql}. La interfaz de línea de comandos recibe un
archivo \texttt{.model}; puede imprimir el SQL en la salida estándar o escribirlo
en un archivo con \texttt{-o}.

\section{Pruebas y resultados}

Los escenarios compartidos de \texttt{casos/} se ejecutan con \texttt{make test}.
El conjunto contiene cinco entradas válidas, que comprueban generación SQL,
traducción de tipos, comentarios, referencias adelantadas y relaciones circulares,
y once entradas inválidas que cubren errores léxicos, sintácticos y semánticos.
Para cada caso se compara la salida o el diagnóstico completo con su archivo de referencia.
La suite también verifica la opción de salida a archivo y errores básicos de uso
y de lectura de archivos.

La ejecución termina correctamente solo cuando todos los escenarios coinciden
con sus resultados esperados. La herramienta no necesita un servidor de base
de datos para las pruebas: estas validan el SQL producido y los diagnósticos,
pero no ejecutan el DDL en una instancia PostgreSQL.

\section{Conclusiones y trabajo futuro}

Transpilador-Model-SQL demuestra una cadena de compilación completa para un
lenguaje específico de dominio acotado: Flex y Bison implementan las etapas de
análisis, el árbol de sintaxis abstracta separa sintaxis de semántica y la etapa
de generación produce DDL de PostgreSQL. La tabla de símbolos y el análisis en dos
pasadas habilitan validaciones de relaciones, incluso cuando hay referencias
adelantadas. El Makefile y los casos compartidos simplifican la compilación y
la demostración del funcionamiento.

Como extensiones posibles se identifican soportar otros dialectos SQL, ampliar
las restricciones disponibles y ejecutar las salidas en PostgreSQL como prueba
adicional de integración.

\begin{thebibliography}{8}

\bibitem{aho}
Aho, A.V., Lam, M.S., Sethi, R., Ullman, J.D.: Compilers: Principles,
Techniques, and Tools, 2nd edn. Pearson (2006)

\bibitem{flex}
The Flex Project: Flex: The Fast Lexical Analyzer. Manual.
\url{https://westes.github.io/flex/manual/}

\bibitem{bison}
Free Software Foundation: GNU Bison Manual.
\url{https://www.gnu.org/software/bison/manual/}

\bibitem{postgresql}
PostgreSQL Global Development Group: PostgreSQL Documentation: Constraints
and Data Types.
\url{https://www.postgresql.org/docs/current/ddl-constraints.html};
\url{https://www.postgresql.org/docs/current/datatype-numeric.html}

\end{thebibliography}

\end{document}
